% Source: https://artofproblemsolving.com/community/c3h2415047_2015_school_sprint_round_29
% Stable identity: aops-topic-2415047-q01
A semicircle of radius $r$ has length $\pi r$ (half of the circumference $2\pi r$). The radii are $1,\dfrac12,\dfrac14,\dots$:

\begin{asy}
usepackage("amsmath");
size(6cm);
defaultpen(fontsize(10pt));
real r=1; real x0=0; real c; bool up=true;
for(int k=0;k<9;++k){
 if(k%2==0){ c=x0+r; draw(arc((c,0),r,180,0),blue+linewidth(0.9)); x0=c+r; }
 else { c=x0-r; draw(arc((c,0),r,0,-180),blue+linewidth(0.9)); x0=c-r; }
 r=r/2; }
draw((-0.1,0)--(2.2,0),gray+dashed);
dot((0,0)); label("start",(0,0),W);
draw((1,0)--(1,1),gray); label("$1$",(1,0.5),W,fontsize(8pt));
draw((1.5,0)--(1.5,-0.5),gray); label("$\dfrac12$",(1.5,-0.25),E,fontsize(8pt));
\end{asy}

So the total length is the sum of a geometric series with first term $\pi$ and ratio $\dfrac12$:
\[\pi\left(1+\dfrac12+\dfrac14+\cdots\right)=\dfrac{\pi}{1-\dfrac12}=\boxed{2\pi}.\]
